\documentclass[]{resume-openfont}
\geometry{hmargin=0.75cm, vmargin=0.5cm}   % slightly tighter than the class default (1cm / 0.94cm) to keep two pages
\renewcommand{\sectionsep}{\vspace{3pt}}    % class default 7pt
% class section style with a tighter line height (class: 16pt/24pt) to keep two pages
\titleformat{\section}{\color{headings}\scshape\customLatoLight\fontsize{16pt}{20pt}\selectfont\raggedright\uppercase}{}{0em}{}
\usepackage[bibstyle=publist,plauthorhandling=highlight,nameorder=given-family,minbibnames=10,maxbibnames=50]{biblatex}
\DeclareDelimFormat[bib]{nametitledelim}{\newline\bibsentence}
\usepackage{academicons}
\usepackage{xcolor}

\plauthorname[Shivam]{Singh}
\pagestyle{fancy}
\resetHeaderAndFooter

\usepackage{tikz}

\newcommand{\ExternalLink}{%
    \tikz[x=1.2ex, y=1.2ex, baseline=-0.05ex]{% 
        \begin{scope}[x=1ex, y=1ex]
            \clip (-0.1,-0.1) 
                --++ (-0, 1.2) 
                --++ (0.6, 0) 
                --++ (0, -0.6) 
                --++ (0, -0.6) 
                --++ (0.6, 0) 
                --++ (0, -1);
            \path[draw, 
                line width = 0.5, 
                rounded corners=0.5] 
                (0,0) rectangle (1,1);
        \end{scope}
        \path[draw, line width = 0.5] (0.5, 0.5) 
            -- (1, 1);
        \path[draw, line width = 0.5] (0.6, 1) 
            -- (1, 1) -- (1, 0.6);
        }
    }

%--------------------------------------------------------------
% Convenience command - make it easy to fill template

% Create job position command. Parameters: company, position, location, when
\newcommand{\resumeHeading}[4]{\runsubsection{\uppercase{#1}}\descript{ | #2}\hfill\location{#3 | #4}\fakeNewLine}

% Create education heading. Parameters: Name, degree, location, when
\newcommand{\educationHeading}[4]{%
\runsubsection{#1}
\descript{| #2 \; (#4)}
\hfill
CGPA: #3
\fakeNewLine
}

% Create project heading. Parameters: Name, link, Tech stack
\newcommand{\projectHeading}[3]{\Project{#1}{#2}
\descript{#3}\\}

\newcommand{\projectHeadingWithDate}[4]{\Project{#1}{#2}
\descript{#3 | #4}\\}

% Parameters: courses
\newcommand{\courseWork}[1]{\textbf{Coursework:} #1}

% Parameters: courses
\newcommand{\teacherAssistant}[1]{\textbf{Teacher Assistant (TA):} #1}
 
%--------------------------------------------------------------
\begin{document}
%--------------------------------------------------------------
%     Profile
%--------------------------------------------------------------
\newcommand{\yourName}{Shivam Singh}

\newcommand{\yourEmail}{shivamsingh3458@gmail.com}
\newcommand{\githubUserName}{sssshivvvv}
\newcommand{\linkedInUserName}{sssshivvvv}
\newcommand{\yourWebsiteLink}{https://shivhere.github.io}
\newcommand{\scholarID}{Wcn55lQAAAAJ}
\newcommand{\phone}{+91 8130870896}


\begin{center}
    \Huge \scshape \latoRegular{\yourName} \\ \vspace{1pt}
    \small
    \href{mailto:\yourEmail}{\faEnvelope \ \underline{\yourEmail}} $|$ 
    \href{https://www.linkedin.com/in/\linkedInUserName}{\faLinkedin \ \underline{LinkedIn}} $|$ 
    \href{https://github.com/\githubUserName}{\faGithub \ \underline{GitHub}} $|$ 
    \href{\yourWebsiteLink}{\faGlobe \ \underline{shivhere.github.io}} $|$ 
    \href{https://scholar.google.com/citations?user=\scholarID}{\aiGoogleScholar \ \underline{Google Scholar}} $|$
    Phone:\phone
\end{center}

\vspace{-4pt}
{\centering
    \textbf{AI Researcher} with \textbf{3+ years of experience} in Vision-Language Models, multimodal reasoning, video understanding, and embodied AI.\par}
\vspace{3pt}

%--------------------------------------------------------------
%     Education
%--------------------------------------------------------------
\section{Education}

\educationHeading
{M.S. -- Electronics and Communication Engineering}
{IIIT Hyderabad}
{8.33/10}
{2023--2026}

% \educationHeading
% {M.S. -- Control and Optimization}
% {IIT Madras}
% {8.50/10}
% {2022}

\educationHeading
{B.Tech -- Electronics and Communication Engineering}
{RKGIT, Ghaziabad}
{8.07/10}
{2017--2021}
\sectionsep


%--------------------------------------------------------------
%     Experience
%--------------------------------------------------------------


\section{Research \& Professional Experience}

\resumeHeading
{pixelstwo}
{Collaboration}
{Remote}
{Sep 2026 - Present}

\textit{\textbf{Focus: Video Data Quality \& Real-to-Sim / Sim-to-Real Systems}}
\sectionsep

\begin{itemize}[itemsep=0pt, parsep=0pt, topsep=1pt]
\item Performing \textbf{quality control of video data}, ensuring recordings are reliable and consistent for downstream training and evaluation.
\item Working on \textbf{real-to-sim and sim-to-real systems}.
\end{itemize}
\sectionsep


\resumeHeading
{intuition.dev}
{Collaboration}
{Remote}
{Aug 2026 - Sep 2026}

\textit{\textbf{Focus: Real-World Video Pipelines}}
\sectionsep

\begin{itemize}[itemsep=0pt, parsep=0pt, topsep=1pt]
\item \textbf{\href{https://shivhere.github.io/assets/data/intuition_work.pdf}{Label-free Segmentation and Grounded Description of Egocentric Work Videos: }}\hfill \textit{\textbf{2026}}
\begin{itemize}
\item Built a \textbf{geometry-first, label-free pipeline} that converts raw body-camera recordings of manual work into \textbf{clean, standalone sub-task clips}: removed inactive/noisy footage with classical detectors and detected \textbf{task-transition points} from metric camera trajectories using purely classical image-processing heuristics, without relying on large models for segmentation and prior annotations.
\item Structured raw videos for downstream \textbf{VLM-based annotation} by isolating relevant task segments before captioning, \textbf{reducing VLM inference cost} and enabling more consistent, task-level annotations under a fixed JSON schema with grounded actions.    [\href{https://shivhere.github.io/assets/data/intuition_work.pdf}{ Technical Report \ExternalLink}]
\end{itemize}
\end{itemize}
\sectionsep


\resumeHeading
{AMD India Private Limited}
{Research Fellow}
{AI Models \& Agents Team}
{Jun 2025 - July 2026}

\textit{\textbf{Focus: VLMs, Video Reasoning \& Spatial Understanding}}
\sectionsep

\textbf{Supervisors:} Saptarshi Majumder, Pratik Prabhanjan Brahma
\sectionsep

\begin{itemize}[itemsep=0pt, parsep=0pt, topsep=1pt]
\item \textbf{Spatial Video Understanding from Raw RGB Videos:}\hfill \textit{\textbf{2026}}
\begin{itemize}
    \item Proposed \textit{SlotMem}, a \textbf{memory-augmented} transformer framework for object tracking, spatial understanding, and long-horizon video reasoning from raw monocular RGB videos.
    
    \item Developed a \textbf{scalable pipeline} combining Depth-Anything-3, YOLO-World, and Vision-Language Models to reconstruct 3D scenes, estimate camera trajectories, and generate supervision for spatial, temporal, and goal-oriented reasoning.
    
    \item Trained SlotMem using our automatically generated supervision together with diverse videos from VLM-3R and VICA-322K, achieving \textbf{state-of-the-art} performance on OST-Bench (\textbf{44.4\%}) and outperforming all prior baselines.
\end{itemize}
\vspace{0.05cm}

\item \textbf{Pause and Think: Video-Grounded Assistive Action Suggestion:}\hfill \textit{\textbf{2025}}

\begin{itemize}
    \item Proposed a reasoning-centric framework for video-grounded assistive action suggestion by introducing \textit{Pause-and-Think-T} (\textbf{$\sim$10k} structured reasoning QA pairs) with explicit reasoning supervision and \textit{Pause-and-Think-B}, a \textbf{300-sample} benchmark for grounded contextual understanding and goal-oriented planning across Epic-Kitchens, Assembly101, and Ego4D.

    \item Built a scalable reasoning-data pipeline and fine-tuned a compact \textbf{4B} VLM on 8$\times$AMD Instinct MI325 GPUs, achieving \textbf{58.0\%} accuracy---matching Qwen3-VL-235B (\textbf{58.9\%}, \textbf{59$\times$ larger}) and GPT-5.2 on scene understanding, outperforming GPT-4o, with strong zero-shot generalization on EgoThink and TempCompass.

    \item Deployed the optimized model on an \textbf{AMD Ryzen\textsuperscript{TM} AI} platform for real-time on-device inference; as \textbf{First Author}, led the research, algorithm design, dataset creation, model training, evaluation, deployment, and manuscript preparation. [\href{https://github.com/sssshivvvv/pause-and-think}{Code \ExternalLink} | \href{https://arxiv.org/abs/2606.00616}{Paper \ExternalLink}]
\end{itemize}

\end{itemize}
\sectionsep


\resumeHeading
{International Institute of Information Technology - Hyderabad}
{M.S. Researcher}
{\href{https://robotics.iiit.ac.in/}{Robotics Research Centre, Hyderabad}}
{Jan 2023 - May 2025}

\textit{\textbf{Focus: Embodied AI \& Robot Planning}}
\sectionsep

\textbf{Direct Supervisor:} \href{https://robotics.iiit.ac.in/faculty_mkrishna/}{Prof. Madhava Krishna}\\
\textbf{Collaborators:} \href{https://homepages.inf.ed.ac.uk/msridhar/}{Dr. Mohan Sridharan}, \href{https://krrish94.github.io/}{Dr. Krishna Murthy Jatavallabhula}, \href{https://sites.google.com/view/brojeshwar/home}{Dr. Brojeshwar Bhowmick}
\sectionsep

\begin{itemize}[itemsep=0pt, parsep=0pt, topsep=1pt]
\item \textbf{Hybrid Planning with LLMs \& Probabilistic Models} (\textit{ECMR 2025}) \hfill \textit{\textbf{2025}}\\
\textbf{Mentor:} Guided framework design and ablation experiments for task planning under uncertainty with probabilistic effects in multi-agent settings. [\href{https://arxiv.org/abs/2602.19518}{Paper \ExternalLink}]
\vspace{0.03cm}

\item \textbf{LLM \& Knowledge Graphs for Task Decomposition} (\textit{IEEE ICRA 2025}) \hfill \textit{\textbf{2025}}\\
\textbf{Project Lead:} Designed a framework that integrates Large Language Models (LLMs), Knowledge Graphs (KGs), and Human-in-the-Loop (HITL) feedback to enable embodied agents to perform long-horizon planning for unseen tasks and \textbf{progressively adapt} to new environments. [\href{https://sssshivvvv.github.io/adaptbot/}{Website \ExternalLink} | \href{https://arxiv.org/pdf/2502.02067}{Paper \ExternalLink}]
\vspace{0.03cm}

\item \textbf{Task Anticipation for Multi-Agent Systems} (\textit{ICRA \& RSS 2024 Workshops}) \hfill \textit{\textbf{2024}}\\
\textbf{Project Lead:} Extended LLM-based task anticipation to long-horizon planning in \textbf{multi-agent settings}, enabling robotic agents to anticipate domain-specific tasks, adapt to unexpected changes in human behavior, and handle environmental uncertainties for effective human--robot collaboration. [\href{https://dataplan-hrc.github.io/}{Website \ExternalLink} | \href{https://arxiv.org/pdf/2404.03587}{Paper \ExternalLink}]
\vspace{0.03cm}

\item \textbf{Task Anticipation using LLMs} (\textit{IEEE ICRA 2024}) \hfill \textit{\textbf{2023}}\\
\textbf{Co-author \& Project Lead:} Developed an intelligent household agent that performs\textbf{ long-horizon planning }by anticipating user actions through learning task execution patterns from user behavior and preferences, leveraging Large Language Models (LLMs) and Planning Domain Definition Language (PDDL)-based task and action representations.[\href{https://raraghavarora.github.io/ahsoka/}{Website \ExternalLink} | \href{https://arxiv.org/pdf/2502.02066}{Paper \ExternalLink}]
\end{itemize}
\sectionsep

\resumeHeading{AlsoEnergy – India}{Technical Support Engineer}{April 2021}{Dec 2021}
\vspace{-0.3em}
\begin{itemize}
\item PowerTrack platform support for commercial solar clients: troubleshooting, data visualization, hardware integration.
\end{itemize}
\sectionsep





\section{Publications \& Workshops}
\begin{itemize}[itemsep=0pt, topsep=1pt]
    \item \href{https://arxiv.org/abs/2606.00616}{\textbf{Pause and Think: A Dataset and Benchmark for Video-Grounded Assistive Action Suggestion: \ExternalLink}} \textbf{Shivam Singh}, S Majumder, P Prabhanjan, Z Liu, E Barsoum. \\
    \textit{Accepted at \href{https://2026.ieee-iros.org/}{IEEE/RSJ International Conference on Intelligent Robots and Systems (IROS) 2026\ExternalLink ~\textbf{(A*)}}. Work done at AMD (Advanced Micro Devices, Inc.).}
    \item \href{https://arxiv.org/abs/2602.19518}{\textbf{Anticipate, Adapt, Act: A Hybrid Framework for Task Planning: \ExternalLink}} N Dash, A Kaura, \textbf{Shivam Singh}, R Singh, S Banerjee, M Sridharan, M Krishna. \\
    \textit{Accepted at IEEE European Conference on Mobile Robots (ECMR) 2025.}
    \item 
    \href{https://sssshivvvv.github.io/adaptbot/}{\textbf{\textcolor{black}{AdaptBot: Combining LLM with Knowledge Graphs and Human Input for Generic-to-Specific Task Decomposition and Knowledge Refinement: \ExternalLink}}} \textbf{Shivam Singh}, K Swaminathan, N Dash, R Singh, S Banerjee, M Sridharan, M Krishna.\\ \textit{Accepted at \href{https://2025.ieee-icra.org/}{IEEE International Conference on Robotics and Automation (ICRA) 2025\ExternalLink~\textbf{(A*)} }, Atlanta, US. 
    {\textbf{\color{black}{\href{https://2025.ieee-icra.org/}{\textcolor{black}}}}}}
    % {\faLink}}}}}
    \item {\textbf{\color{black}{\href{https://openreview.net/forum?id=XmInDYtUPa}{\textcolor{black}{Anticipate \& Collab: Data-driven Task Anticipation and Knowledge-driven Planning for Human-robot Collaboration \ExternalLink }}}}}: \textbf{Shivam Singh$^\ast$}, K Swaminathan$^\ast$, R Arora$^\ast$, D Das, S Banerjee, M Sridharan, M Krishna. \textit{Accepted in \href{https://sites.google.com/view/rss-taskspec}{RSS2024 Workshop on Task Specification for General-Purpose Intelligent Robots \ExternalLink}, Delft, Netherlands}
    \item \href{https://events.infovaya.com/uploads/documents/pdfviewer/37/ff/133990-3210.pdf}{\textbf{Anticipate \& Act: Integrating LLMs and Classical Planning for Efficient Task Execution in Household Environments: \ExternalLink}} R Arora*, \textbf{Shivam Singh*},  K Swaminathan, A Datta, S Banerjee, B Bhowmick, KM Jatavallabhula, M Sridharan, M Krishna. \\
    \textit{Accepted at \href{https://2024.ieee-icra.org/}{IEEE International Conference on Robotics and Automation (ICRA) 2024 \ExternalLink ~\textbf{(A*)}}, Yokohama, Japan.}
    \item \href{https://drive.google.com/file/d/1sHZiiRoblnyshA-U7Gh8xrvXestDAaEH/view}{\textbf{DaTAPlan: Data-driven Task Anticipation and Knowledge-driven Planning for Human-robot Collaboration: \ExternalLink }} \textbf{Shivam Singh*},  K Swaminathan*, R Arora*, R Singh, A Datta, D Das, S Banerjee, M Sridharan, and M Krishna. \\
    \textit{Accepted in ICRA 2024 Workshop on \href{https://sites.google.com/view/icra2024cookingrobotics/}{Cooking Robotics: Perception and motion planning \ExternalLink}, Yokohama, Japan.}
\end{itemize}
\setlength{\linewidth}{18cm}

\section{Patents}
\begin{itemize}[itemsep=0pt, topsep=1pt]
    \item \href{https://patents.google.com/patent/US20260216870A1/en}{\textbf{Method for Automatic Knowledge Refinement and Task Execution \ExternalLink}}: \textbf{Shivam Singh}, S Banerjee, K Swaminathan, N Dash, R Singh, M Sridharan, and M Krishna. \textit{US Patent App. 19420895}, published July 30, 2026.
    \item \href{https://patents.google.com/patent/US20250342322A1/en}{\textbf{Method and System for Task Anticipation by Integrating Large Language Models and Classical Planning \ExternalLink}}: \textbf{Shivam Singh}, S Banerjee, B Bhowmick, R Arora, A Datta, K Swaminathan, M Krishna, M Sridharan, and KM Jatavallabhula. \textit{US Patent App. 19089689}, published November 6, 2025.
\end{itemize}



\section{Technical Skills}

\textbf{Research Areas:}
LLMs, VLMs, VLAs, Reasoning Models, Computer Vision, Multimodal Learning, Agentic AI, Embodied AI

\textbf{Training \& Inference:}
PyTorch, Hugging Face, LLaMA-Factory, SFT, RLHF, LoRA, DeepSpeed, vLLM, ROCm

\textbf{Data \& Infrastructure:}
Multimodal Data Curation, Synthetic Data, Benchmarking, Docker, Slurm, OpenCV, Open3D

\textbf{Programming \& Robotics:}
Python, C++, C, PDDL, RDDL, ROS, Gazebo, AI2-THOR


\end{document}


































































